% Modèle matOS : feuille de travaux dirigés. Compiler avec pdflatex.
\documentclass[11pt,a4paper]{article}
\usepackage[T1]{fontenc}
\usepackage[utf8]{inputenc}
\usepackage[french]{babel}
\usepackage{lmodern,microtype}
\usepackage[margin=2.2cm]{geometry}
\usepackage{amsmath,amssymb,mathtools}
\usepackage{enumitem}
\usepackage{fancyhdr}

\newcounter{exercice}
\newenvironment{exercice}[1][]{%
  \refstepcounter{exercice}\par\medskip\noindent
  \textbf{Exercice \theexercice}\ifx&#1&\else\ (#1)\fi\textbf{.}\ }{\par}
\newcommand{\R}{\mathbb{R}}
\newcommand{\N}{\mathbb{N}}

\pagestyle{fancy}
\lhead{Mention Mathématiques --- L1}
\rhead{Analyse 1 --- 2026--2027}
\cfoot{\thepage}

\begin{document}
\begin{center}
  {\Large\bfseries Feuille de TD n\textsuperscript{o} 1 : suites réelles}
\end{center}

\begin{exercice}[limites]
Étudier la convergence des suites suivantes :
\begin{enumerate}[label=\alph*)]
  \item $u_n = \dfrac{3n^2 + 1}{n^2 + n}$ ;
  \item $v_n = \sqrt{n+1} - \sqrt{n}$ ;
  \item $w_n = \left(1 + \dfrac{1}{n}\right)^{n}$.
\end{enumerate}
\end{exercice}

\begin{exercice}
Soit $(u_n)$ une suite croissante et majorée. Montrer qu'elle converge vers
$\sup \{ u_n : n \in \N \}$.
\end{exercice}

\begin{exercice}[à chercher avant la séance]
Montrer que toute suite convergente est bornée. La réciproque est-elle vraie ?
\end{exercice}

\end{document}
